\chapter{Literature Survey}

This chapter looks at the work that other researchers have already done in the areas this
project touches. It covers wearable EEG devices, publicly available motor imagery
datasets, the methods used to extract features and classify EEG signals, and the use of
EEG to measure a person's mental and cognitive state. A short section on decoding speech
from EEG is also included, since that is the direction this work is expected to move in
later. The aim of this chapter is to understand what has already been done, what problems
are still open, and where this project fits in.

%------------------------------------------------------------------
\section{Wearable EEG Systems}
%------------------------------------------------------------------

EEG devices have come a long way from the large machines used in hospitals. Today there
are small, portable EEG headsets that a person can wear while going about normal
activities. This matters for any brain-computer interface (BCI) that is meant to be used
outside a lab.

Mullen and colleagues put together a complete wearable system with dry electrodes and
showed that it could follow brain activity and mental state live, while the person was
wearing it \cite{mullen2015wearable}. Because no gel is used, the headset is quick to put
on and stays comfortable for a long time. Their recordings were good enough for practical
use, although still a little noisier than a gel cap in a laboratory.

Kappel and colleagues went a step further and moved the dry electrodes into the ear
canal, so that the recorder looks like an ordinary earphone \cite{kappel2019earEEG}. The
signals they picked up were weaker than with electrodes on the scalp, but the device was
so unobtrusive that it could be worn through a normal day, which is the kind of comfort a
practical brain-computer interface needs.

Both of these studies are directly relevant to the hardware used in this project. The GELO
BCI headset used here is a dry-electrode headphone with only four EEG channels (A1, A2, C3
and C4) and a single reference at Cz, sampled at 250 Hz. It gains the comfort and quick
setup reported in these studies, but it also inherits their limitations: fewer channels,
weaker signals and more noise than a research-grade cap. This trade-off shapes much of the
work reported in the later chapters.

%------------------------------------------------------------------
\section{Motor Imagery EEG Datasets}
\label{sec:mi_datasets}
%------------------------------------------------------------------

Motor imagery (MI) means imagining a movement without actually making it. When a person
imagines moving the left or right hand, the activity over the motor cortex changes, and
these changes can be picked up by electrodes placed at C3 and C4. Building a classifier for
this needs a good amount of labelled data, which is hard for a student project to collect
from scratch.

The first reference dataset used in this project comes from Liu and colleagues, who
released motor imagery recordings from people in the first month after a stroke
\cite{liu2024mi}. Fifty patients took part, the youngest aged 31 and the oldest 77. Every
patient did 40 trials, half imagining gripping with the left hand and half with the right,
which adds up to 2000 trials. A trial lasted 8 s: 2 s of instruction, 4 s of imagery and
2 s of rest. The recordings were made at 500 Hz with a portable wireless cap using saline
instead of gel, with 30 EEG electrodes, 2 EOG electrodes and a marker channel, referenced
to CPz. The authors share both the raw and the cleaned data, together with clinical notes
on each patient.

This dataset is important for the present work for two reasons. First, it supplies a
labelled, well-documented set of left and right MI trials that can be used to train a
classifier, which is the first objective of this project. Second, the paper reports its own
analysis of the data, including event-related desynchronisation and a set of classification
results, which gives a clear baseline to compare against. The channel layout is also
convenient, as C3 and C4 are available in the dataset and are the same two positions
recorded by the headset used here. The main limitation is that the subjects are stroke
patients recorded with a 30-channel system, while the data collected in this project comes
from healthy subjects using only four channels. This difference between the training data
and the test data turns out to matter a great deal, as the results chapter shows.

A second public dataset was used to fill this gap. Dataset 2b of BCI Competition IV, which
was released for the competition in 2008 and is called the BCI 2008 dataset in this report,
was recorded at Graz University of Technology~\cite{leeb2007virtual,
tangermann2012review}. It contains nine healthy, right-handed people who imagined moving
the left or the right hand whenever an arrow pointing to that side appeared on a screen.
Each person came back for five sessions on different days. The first two sessions were
plain cued imagery without feedback, about 120 trials each, and in the last three, about
160 trials each, a smiley face on the screen moved left or right according to the
person's own brain signal. Only three EEG channels were recorded, over C3, Cz and C4, each
as a bipolar pair of closely spaced electrodes, at 250 Hz with a 0.5--100 Hz band-pass
filter and a 50 Hz notch filter. Three EOG channels were also recorded, but they are meant
only for spotting eye movements and not for classification. The dataset has 6520 trials
in total, half for each hand.

This dataset suits the present work for three reasons. The subjects are healthy, as in
the recordings made for this project. The sampling rate is the same 250 Hz as the GELO
headset, and the layout is almost the same, with the two motor cortex electrodes at C3 and
C4 plus one at Cz. Finally, it comes with a fixed split, three sessions for training and
two later sessions for testing, so a classifier can be checked on days it has never seen.
Many methods have been tested on it, and the winning entry of the competition used filter
bank CSP~\cite{ang2012fbcsp}, which is one of the methods tried in this project.

%------------------------------------------------------------------
\section{Feature Extraction and Classification for Motor Imagery}
%------------------------------------------------------------------

Once the data is available, the next problem is turning the raw signal into features and
training a model on them. This is harder than it sounds, because EEG signals vary a lot
between people and even between sessions for the same person.

The standard way of looking at motor imagery is through event-related desynchronisation
(ERD) and event-related synchronisation (ERS). ERD is a drop in power in a given frequency
band during a task and ERS is a rise in power. For motor imagery, power in the alpha
(8--15 Hz) and beta (15--30 Hz) bands usually drops over the motor cortex on the side
opposite to the imagined hand \cite{pfurtscheller1999erd}. The reference dataset paper computes this by filtering C3
and C4 between 8 and 30 Hz, squaring the signal to get power, averaging across trials and
comparing against a baseline taken just before the task \cite{liu2024mi}. The same paper
also benchmarks several classifiers on its own data for the left versus right MI problem,
using the 8--30 Hz band and a 4 second window after the cue. Common Spatial Pattern with
Linear Discriminant Analysis (CSP+LDA) \cite{ramoser2000csp} reached about 55.6\%, Filter
Bank CSP with a Support Vector Machine (FBCSP+SVM) \cite{ang2012fbcsp, cortes1995svm}
about 57.6\%, a Riemannian approach \cite{barachant2012riemann} about 61.2\%, and the best
method reported, a temporal window filter bank combined with a Riemannian mean classifier,
reached 72.21\%. These numbers are a useful reminder that even with 30 channels and clean
laboratory recordings, binary motor imagery classification is a difficult problem, and that
accuracies in the 50--70\% range are normal rather than poor.

Beyond CSP and SVM, deep networks are the other big family of methods. A survey by Xu
and Xia walks through the network types that have been tried on EEG and the usual ways of
preparing the signal for them \cite{xu2023review}. Its main warning is about data size:
deep networks learn well only from large amounts of data, and most EEG datasets are small.
For a project that records its own data from a handful of people, a simple model like an
SVM is therefore the sensible first choice.

EEG Conformer, proposed by Song and colleagues, mixes two ideas
\cite{song2023conformer}. Convolution layers look for short local patterns in the signal,
and an attention stage then relates those patterns to each other over the whole trial. The
authors also show maps of which parts of the input the network relied on, and they report
higher accuracy than CSP on the usual motor imagery benchmarks. It is a good candidate for
later work, once there is enough data and computing power to train it.

Panachakel and Ramakrishnan used transfer learning for imagined speech
\cite{jerrin2021imagined}. Their network learns first from many people and is then tuned to
a new person with only a little of that person's data, which shortens the setup for each
new user. Accuracy still varied clearly from one person to the next, showing that
differences between people remain a hard problem. The
same problem appears in this project, where a model trained on one dataset is tested on
data recorded from a different group of people with a different device.

Su and colleagues built a brain-computer interface for hand rehabilitation after stroke
that uses two kinds of brain response at once: motor imagery, and the steady response the
visual cortex gives to a flickering light \cite{su2024hybrid}. Their model keeps updating itself as the patient's signals change over time. Using both signals gave
higher accuracy than either alone, and patients' hand function improved, which shows that
an EEG classifier can be of real clinical use.

%------------------------------------------------------------------
\section{Decoding Speech from EEG}
%------------------------------------------------------------------

A related line of research tries to work out what a person is saying, or only thinking of
saying, from EEG alone. This is not part of the current objectives but it is the direction
this project is expected to extend into, so the main findings are worth noting.

Krishna and colleagues recorded EEG while people spoke and found that a speech recogniser
could use it, both when the audio was noisy and when there was no audio at all
\cite{krishna2019nospeech}. This mattered because it meant
EEG on its own carries enough information about speech intent to be used for recognition.
The experiments were done in a controlled lab, so how well this holds in noisy everyday
conditions is still not known. In a follow-up study the same group tried to produce speech sound directly from EEG and
managed a partial reconstruction of the spoken words \cite{krishna2020eeg}. The method depended on specific tasks and had to be trained
separately for each person, which makes it hard to use on a new user without starting over.

EEG recorded during both imagined and spoken speech has also been used to classify the
types of sounds being thought about or said \cite{zhao2015phonological}. This confirmed
that imagined speech, where no sound is made at all, still leaves a pattern in EEG that can
be detected, which is the basic assumption behind any silent communication system. On the
listening side, the auditory EEG decoding challenge for ICASSP 2024 focused on working out
what a person is listening to from their EEG \cite{bollens2025auditory}. The dataset and the
evaluation method were released publicly, which helps the wider research community compare
results fairly.

%------------------------------------------------------------------
\section{EEG for Cognitive and Mental State Assessment}
%------------------------------------------------------------------

Apart from decoding what a person intends to do or say, EEG is also widely used to measure
how a person's mind is working at a given moment. This is the third objective of this
project, which is to compute a cognitive score from EEG frequency bands.

Bird and colleagues used a low-cost headband and a set of statistical features to sort
recordings into relaxed, neutral and concentrating states with machine learning
\cite{bird2018mental}, and found that these states can be told apart with good accuracy. The work looked at states in isolation, however,
and did not follow how a changing cognitive state affects performance on an actual task.

Jawed and colleagues trained recurrent networks on unprocessed EEG to pick out people who
learn best from visual material \cite{jawed2024visual}. With 34 participants recorded at
rest and while learning, their best model, which combined an LSTM with convolution layers,
was right about 94\% of the time. What is useful here is that the models worked on raw EEG without heavy offline
cleaning, which shows that real-time assessment from a continuous EEG stream is practical.

For driver drowsiness, Ming and colleagues used deep Q-learning, a reinforcement learning
method that learns from feedback instead of from fixed class labels
\cite{ming2021drowsiness}. This is a useful idea because drowsiness builds up gradually and does not
split neatly into classes. A more conventional system for the same problem used spectral
entropy and power spectral density as features, which describe how the brain's electrical
activity is spread across frequency bands, and fed them to boosting classifiers to detect
signs of fatigue \cite{jiri2025drowsiness}. This second approach is closest to the method
used in this project, where relative power in the delta, theta, alpha, beta and gamma bands
is computed from the PSD and then combined into scores. In particular, the stress score used
here is based on the ratio of relative high-beta power to relative alpha power, and the
fatigue score on the ratio of combined theta and alpha power to beta power, both of which
rest on the same band-power relationships these studies rely on. The limitation of such
systems is that they are tuned for one specific state, so the thresholds and reference
values have to be worked out again for any new measure.

%------------------------------------------------------------------
\section{Summary of the Literature Survey}
%------------------------------------------------------------------

The studies reviewed above map onto the three objectives of this project as follows.

\textbf{Building an EEG database for motor imagery.} Wearable dry-electrode systems,
including ear-mounted ones, can record usable EEG with a small number of channels and very
little setup \cite{mullen2015wearable, kappel2019earEEG}, and the open stroke-patient motor
imagery dataset provides a well-documented left and right hand MI recording protocol that
can be followed when collecting new data \cite{liu2024mi}. The BCI Competition IV 2b
dataset adds healthy subjects recorded with only three channels at 250 Hz, which is much
closer to a small wearable headset \cite{leeb2007virtual}.

\textbf{Designing a motor imagery classifier.} ERD and ERS analysis of C3 and C4 in the
8--30 Hz band is the accepted way to look at motor imagery, and CSP, filter bank CSP,
LDA, SVM and Riemannian methods form the standard set of classifiers, with reported
accuracies for left versus right MI ranging from about 55\% to 72\%
\cite{liu2024mi}. Deep learning offers further gains \cite{xu2023review,
song2023conformer}, transfer learning helps when moving between users
\cite{jerrin2021imagined}, and adaptive models help when a user's signals drift over time
\cite{su2024hybrid}.

\textbf{Computing a cognitive score from frequency bands.} Relative power in the EEG
frequency bands, together with PSD and entropy features, can be used to separate mental
states and to track attention and fatigue \cite{bird2018mental, jawed2024visual,
ming2021drowsiness, jiri2025drowsiness}.

Three gaps stand out from this review. First, almost all of the motor imagery results
reported above come from research-grade systems with 20 or more channels, while a practical
wearable device may have only a few; how much accuracy is lost in that reduction is not
well documented. Second, models are usually trained and tested on the same dataset and the
same recording device, so very little is known about what happens when a model trained on
one dataset is applied to data recorded elsewhere with different hardware. Third, work on
motor imagery decoding and work on cognitive state monitoring have largely been carried out
separately, even though both are computed from the same recording. This project addresses
these three points by recording left and right motor imagery on a four-channel dry
headset, training classifiers on the public dataset and testing them on the recorded data,
and computing a cognitive score from the same signal within a single system.